\documentclass[10pt,a4paper]{amsart}
\usepackage{amsmath,amssymb,amsthm,graphicx}
% Retained preamble for the Stacks Project chapter {CH} (shared project preamble).
% The chapter file pinned-original-source.tex only carries \input{{preamble}} on its first line,
% so no preamble text lives inside the chapter itself.  The Stacks Project uses one single shared
% preamble for every chapter: the snapshot stages it three times (sources/stacks-chow/preamble.tex,
% sources/stacks-etale-cohomology/preamble.tex and focused-seeds/AGP/sources/stacks-snapshot-20261011/
% preamble.tex) and those copies are identical in size.  The block below is copied mechanically,
% line for line, from sources/stacks-chow/preamble.tex lines 14-18 and 154-217: the Xy-pic setup and
% the native theorem declarations together with the mathematical macros.  Nothing was added, removed
% or redefined; the wrapper class and the amsmath/amssymb/amsthm/graphicx base packages are supplied
% by the assembler.
% For commutative diagrams we use Xy-pic
% Chapter: stacks-spaces-more-morphisms.tex -- More on Morphisms of Spaces.
\usepackage[all]{xy}

% We use 2cell for 2-commutative diagrams.
\xyoption{2cell}

% Theorem environments.
%
\theoremstyle{plain}
\newtheorem{theorem}[subsection]{Theorem}
\newtheorem{proposition}[subsection]{Proposition}
\newtheorem{lemma}[subsection]{Lemma}

\theoremstyle{definition}
\newtheorem{definition}[subsection]{Definition}
\newtheorem{example}[subsection]{Example}
\newtheorem{exercise}[subsection]{Exercise}
\newtheorem{situation}[subsection]{Situation}

\theoremstyle{remark}
\newtheorem{remark}[subsection]{Remark}
\newtheorem{remarks}[subsection]{Remarks}

\numberwithin{equation}{subsection}

% Macros
%
\def\lim{\mathop{\mathrm{lim}}\nolimits}
\def\colim{\mathop{\mathrm{colim}}\nolimits}
\def\Spec{\mathop{\mathrm{Spec}}}
\def\Hom{\mathop{\mathrm{Hom}}\nolimits}
\def\Ext{\mathop{\mathrm{Ext}}\nolimits}
\def\SheafHom{\mathop{\mathcal{H}\!\mathit{om}}\nolimits}
\def\SheafExt{\mathop{\mathcal{E}\!\mathit{xt}}\nolimits}
\def\Sch{\mathit{Sch}}
\def\Mor{\mathop{\mathrm{Mor}}\nolimits}
\def\Ob{\mathop{\mathrm{Ob}}\nolimits}
\def\Sh{\mathop{\mathit{Sh}}\nolimits}
\def\NL{\mathop{N\!L}\nolimits}
\def\CH{\mathop{\mathrm{CH}}\nolimits}
\def\proetale{{pro\text{-}\acute{e}tale}}
\def\etale{{\acute{e}tale}}
\def\QCoh{\mathit{QCoh}}
\def\Ker{\mathop{\mathrm{Ker}}}
\def\Im{\mathop{\mathrm{Im}}}
\def\Coker{\mathop{\mathrm{Coker}}}
\def\Coim{\mathop{\mathrm{Coim}}}

% Boxtimes
%
\DeclareMathSymbol{\boxtimes}{\mathbin}{AMSa}{"02}

%
% Macros for moduli stacks/spaces
%
\def\QCohstack{\mathcal{QC}\!\mathit{oh}}
\def\Cohstack{\mathcal{C}\!\mathit{oh}}
\def\Spacesstack{\mathcal{S}\!\mathit{paces}}
\def\Quotfunctor{\mathrm{Quot}}
\def\Hilbfunctor{\mathrm{Hilb}}
\def\Curvesstack{\mathcal{C}\!\mathit{urves}}
\def\Polarizedstack{\mathcal{P}\!\mathit{olarized}}
\def\Complexesstack{\mathcal{C}\!\mathit{omplexes}}
% \Pic is the operator that assigns to X its picard group, usage \Pic(X)
% \Picardstack_{X/B} denotes the Picard stack of X over B
% \Picardfunctor_{X/B} denotes the Picard functor of X over B
\def\Pic{\mathop{\mathrm{Pic}}\nolimits}
\def\Picardstack{\mathcal{P}\!\mathit{ic}}
\def\Picardfunctor{\mathrm{Pic}}
\def\Deformationcategory{\mathcal{D}\!\mathit{ef}}
\newenvironment{reference}{}{}
\newenvironment{slogan}{}{}
\newenvironment{history}{}{}
\usepackage{cleveref}
\crefname{originallemma}{original lemma}{original lemmas}
\Crefname{originallemma}{Original lemma}{Original lemmas}
\makeatletter
\newlabel{lemma-etale-splits-off-quasi-finite-part}{{02LP}{source}{}{sourcepointer-1}{}}
\newlabel{lemma-etale-splits-off-quasi-finite-part@cref}{{[originallemma][0][]02LP}{[1][0][]source}{}{}{}}
\makeatother
\crefname{originallemma}{original lemma}{original lemmas}
\Crefname{originallemma}{Original lemma}{Original lemmas}
\makeatletter
\newlabel{spaces-morphisms-lemma-finite-proper}{{04NZ}{source}{}{sourcepointer-2}{}}
\newlabel{spaces-morphisms-lemma-finite-proper@cref}{{[originallemma][0][]04NZ}{[1][0][]source}{}{}{}}
\makeatother
\crefname{originallemma}{original lemma}{original lemmas}
\Crefname{originallemma}{Original lemma}{Original lemmas}
\makeatletter
\newlabel{spaces-morphisms-lemma-finite-quasi-finite}{{04NY}{source}{}{sourcepointer-3}{}}
\newlabel{spaces-morphisms-lemma-finite-quasi-finite@cref}{{[originallemma][0][]04NY}{[1][0][]source}{}{}{}}
\makeatother
\crefname{originallemma}{original lemma}{original lemmas}
\Crefname{originallemma}{Original lemma}{Original lemmas}
\makeatletter
\newlabel{spaces-morphisms-lemma-integral-local}{{03ZQ}{source}{}{sourcepointer-4}{}}
\newlabel{spaces-morphisms-lemma-integral-local@cref}{{[originallemma][0][]03ZQ}{[1][0][]source}{}{}{}}
\makeatother
\crefname{originallemma}{original lemma}{original lemmas}
\Crefname{originallemma}{Original lemma}{Original lemmas}
\makeatletter
\newlabel{spaces-morphisms-lemma-locally-quasi-finite}{{06RW}{source}{}{sourcepointer-5}{}}
\newlabel{spaces-morphisms-lemma-locally-quasi-finite@cref}{{[originallemma][0][]06RW}{[1][0][]source}{}{}{}}
\makeatother
\crefname{originallemma}{original lemma}{original lemmas}
\Crefname{originallemma}{Original lemma}{Original lemmas}
\makeatletter
\newlabel{spaces-morphisms-lemma-quasi-finite-at-point}{{06LS}{source}{}{sourcepointer-6}{}}
\newlabel{spaces-morphisms-lemma-quasi-finite-at-point@cref}{{[originallemma][0][]06LS}{[1][0][]source}{}{}{}}
\makeatother
\crefname{originallemma}{original lemma}{original lemmas}
\Crefname{originallemma}{Original lemma}{Original lemmas}
\makeatletter
\newlabel{spaces-morphisms-lemma-universally-closed-quasi-compact}{{04XW}{source}{}{sourcepointer-7}{}}
\newlabel{spaces-morphisms-lemma-universally-closed-quasi-compact@cref}{{[originallemma][0][]04XW}{[1][0][]source}{}{}{}}
\makeatother
\crefname{originallemma}{original lemma}{original lemmas}
\Crefname{originallemma}{Original lemma}{Original lemmas}
\makeatletter
\newlabel{spaces-properties-lemma-points-cartesian}{{03H4}{source}{}{sourcepointer-8}{}}
\newlabel{spaces-properties-lemma-points-cartesian@cref}{{[originallemma][0][]03H4}{[1][0][]source}{}{}{}}
\makeatother
\title{Four equivalent guises of finiteness for morphisms of algebraic spaces}
\author{The Stacks Project Authors}
\date{stacks-project:spaces-more-morphisms.tex:local-snapshot-20261011; GNU FDL 1.2 or later}
\begin{document}
\maketitle
\noindent\textit{Editorial note.} Statement span is lines 7847-7860 of sources/stacks-spaces-more-morphisms/pinned-original-source.tex; proof span is lines 7862-7900, the complete author proof (37 body lines).\par\medskip
\noindent\textit{Editorial note.} Navigation stubs: 8 cited-but-not-included labels are supplied as non-load-bearing pointers with their pinned tags where the snapshot tag index has them.\par\medskip
\noindent\textit{Editorial note.} extra\_packages is empty on purpose: the retained preamble already loads xy with options and a second declaration clashes in the assembler wrapper.\par\medskip
\noindent\textit{Editorial note.} Text is retained verbatim, including the author's own typos and placeholders.\par\medskip
\noindent\textit{Editorial note.} Counted B side (both complete); spaces proof materially uses the chapter's etale-splitting machinery.\par\medskip
\noindent\textit{Allowed standard prerequisites (editorial).}\begin{itemize}
\item Explicit prerequisite boundary. Morphisms of Spaces: finite-proper, finite-quasi-finite, locally-quasi-finite, universally-closed-quasi-compact, quasi-finite-at-point, integral-local, image and point lemmas of Properties of Spaces; chapter lemma etale-splits-off-quasi-finite-part.
\end{itemize}\medskip
\ifdefined\hypertarget\hypertarget{sourcepointer-1}{}\fi
\noindent\textit{Source reference.} Original lemma 02LP: sources/stacks-spaces-more-morphisms/pinned-original-source.tex, source lines 7472--7482. Cited by the retained statement or proof but not part of the collected spans; supplied only as a navigation pointer. The argument does not depend on its content.\par\medskip
\ifdefined\hypertarget\hypertarget{sourcepointer-2}{}\fi
\noindent\textit{Source reference.} Original lemma 04NZ: sources/stacks-spaces-more-morphisms/pinned-original-source.tex, source lines 7900--7900. Cited by the retained statement or proof but not part of the collected spans; supplied only as a navigation pointer. The argument does not depend on its content.\par\medskip
\ifdefined\hypertarget\hypertarget{sourcepointer-3}{}\fi
\noindent\textit{Source reference.} Original lemma 04NY: sources/stacks-spaces-more-morphisms/pinned-original-source.tex, source lines 7900--7900. Cited by the retained statement or proof but not part of the collected spans; supplied only as a navigation pointer. The argument does not depend on its content.\par\medskip
\ifdefined\hypertarget\hypertarget{sourcepointer-4}{}\fi
\noindent\textit{Source reference.} Original lemma 03ZQ: sources/stacks-spaces-more-morphisms/pinned-original-source.tex, source lines 7900--7900. Cited by the retained statement or proof but not part of the collected spans; supplied only as a navigation pointer. The argument does not depend on its content.\par\medskip
\ifdefined\hypertarget\hypertarget{sourcepointer-5}{}\fi
\noindent\textit{Source reference.} Original lemma 06RW: sources/stacks-spaces-more-morphisms/pinned-original-source.tex, source lines 7900--7900. Cited by the retained statement or proof but not part of the collected spans; supplied only as a navigation pointer. The argument does not depend on its content.\par\medskip
\ifdefined\hypertarget\hypertarget{sourcepointer-6}{}\fi
\noindent\textit{Source reference.} Original lemma 06LS: sources/stacks-spaces-more-morphisms/pinned-original-source.tex, source lines 7900--7900. Cited by the retained statement or proof but not part of the collected spans; supplied only as a navigation pointer. The argument does not depend on its content.\par\medskip
\ifdefined\hypertarget\hypertarget{sourcepointer-7}{}\fi
\noindent\textit{Source reference.} Original lemma 04XW: sources/stacks-spaces-more-morphisms/pinned-original-source.tex, source lines 7900--7900. Cited by the retained statement or proof but not part of the collected spans; supplied only as a navigation pointer. The argument does not depend on its content.\par\medskip
\ifdefined\hypertarget\hypertarget{sourcepointer-8}{}\fi
\noindent\textit{Source reference.} Original lemma 03H4: sources/stacks-spaces-more-morphisms/pinned-original-source.tex, source lines 7900--7900. Cited by the retained statement or proof but not part of the collected spans; supplied only as a navigation pointer. The argument does not depend on its content.\par\medskip
\begin{lemma}
\label{lemma-characterize-finite}
Let $S$ be a scheme. Let $f : X \to Y$ be a morphism of algebraic spaces
over $S$. The following are equivalent:
\begin{enumerate}
\item $f$ is finite,
\item $f$ is proper and locally quasi-finite,
\item $f$ is proper and $|X_k|$ is a discrete space for every morphism
$\Spec(k) \to Y$ where $k$ is a field,
\item $f$ is universally closed, separated, locally of finite type
and $|X_k|$ is a discrete space for every morphism $\Spec(k) \to Y$
where $k$ is a field.
\end{enumerate}
\end{lemma}

\begin{proof}
We have (1) $\Rightarrow$ (2) by
Morphisms of Spaces, Lemmas \ref{spaces-morphisms-lemma-finite-proper},
\ref{spaces-morphisms-lemma-finite-quasi-finite}.
We have (2) $\Rightarrow$ (3) by
Morphisms of Spaces, Lemma
\ref{spaces-morphisms-lemma-locally-quasi-finite}.
By definition (3) implies (4).

\medskip\noindent
Assume (4). Since $f$ is universally closed it is quasi-compact
(Morphisms of Spaces, Lemma
\ref{spaces-morphisms-lemma-universally-closed-quasi-compact}).
Pick a point $y$ of $|Y|$. We represent $y$ by a
morphism $\Spec(k) \to Y$. Note that $|X_k|$ is finite discrete
as a quasi-compact discrete space. The map $|X_k| \to |X|$ surjects
onto the fibre of $|X| \to |Y|$ over $y$
(Properties of Spaces, Lemma \ref{spaces-properties-lemma-points-cartesian}).
By
Morphisms of Spaces, Lemma \ref{spaces-morphisms-lemma-quasi-finite-at-point}
we see that $X \to Y$ is quasi-finite at all the points of the fibre
of $|X| \to |Y|$ over $y$.
Choose an elementary \'etale neighbourhood $(U, u) \to (Y, y)$
and decomposition $X_U = V \amalg W$ as in
Lemma \ref{lemma-etale-splits-off-quasi-finite-part}
adapted to all the points of $|X|$ lying over $y$.
Note that $W_u = \emptyset$ because we used all the points
in the fibre of $|X| \to |Y|$ over $y$.
Since $f$ is universally closed we see that
the image of $|W|$ in $|U|$ is a closed set not containing $u$.
After shrinking $U$ we may assume that $W = \emptyset$.
In other words we see that $X_U = V$ is finite over $U$.
Since $y \in |Y|$ was arbitrary
this means there exists a family $\{U_i \to Y\}$
of \'etale morphisms whose images cover $Y$ such that
the base changes $X_{U_i} \to U_i$ are finite.
We conclude that $f$ is finite by
Morphisms of Spaces, Lemma \ref{spaces-morphisms-lemma-integral-local}.
\end{proof}

\end{document}
